\documentclass[11pt]{article}
\usepackage[T1]{fontenc}
\usepackage[utf8]{inputenc}
\usepackage[polish]{babel}
\usepackage{amsmath}
\usepackage{amssymb}
\usepackage{booktabs}
\usepackage{enumitem}
\usepackage[hidelinks]{hyperref}
\usepackage{xcolor}
\usepackage{tikz}
\usepackage{pgfplots}
\pgfplotsset{compat=1.18}
\usetikzlibrary{arrows.meta}

\newcommand{\code}[1]{\texttt{#1}}
\newcommand{\inn}{\ensuremath{\mathrm{in}}}

\definecolor{mkbg}{RGB}{44,110,180}
\definecolor{anbg}{RGB}{160,90,30}

\newcommand{\pyt}[1]{\medskip\noindent\textbf{\textcolor{mkbg}{Pytanie:} #1}\par\smallskip}
\newcommand{\pytan}[1]{\medskip\noindent\textbf{\textcolor{anbg}{Pytanie:} #1}\par\smallskip}

\title{Podsumowania lingwistyczne --- notatka egzaminacyjna\\[3pt]
\large Odpowiedzi do zagadnień (MK --- szczegółowo, AN --- skrótowo)}
\author{}
\date{}

\begin{document}
\maketitle

\noindent\textit{Notatka odpowiada na zebrane zagadnienia egzaminacyjne. Część~\textbf{MK}
(nasz prowadzący) jest rozwinięta dokładnie, część~\textbf{AN} skrótowo, ale z kompletem
pojęć. Każda odpowiedź jest osadzona w faktycznej implementacji projektu --- w nawiasach
podaję klasy/pliki, w których dane pojęcie żyje w kodzie. Domena przykładów: samochody
(moc silnika, masa, długość, spalanie).}

\tableofcontents

% =====================================================================
\section*{Mini-słownik (żeby wszystko miało wspólny język)}
\addcontentsline{toc}{section}{Mini-słownik}

\begin{itemize}[leftmargin=*]
\item \textbf{Przestrzeń rozważań} $X$ --- zbiór wszystkich możliwych wartości cechy
  (np. moc $\in [0, 600]$~KM). W kodzie \code{Universe} (próbkowana siatką punktów).
\item \textbf{Zbiór rozmyty} $A$ --- przypisuje każdemu $x\in X$ stopień przynależności
  $\mu_A(x)\in[0,1]$ (zamiast twardego 0/1). W kodzie \code{FuzzySet} + funkcja
  przynależności \code{MembershipFunction}.
\item \textbf{Etykieta} --- słowo języka naturalnego (np. ,,mocny'') sklejone z jednym
  zbiorem rozmytym (\code{Label}).
\item \textbf{Zmienna lingwistyczna} --- cecha (np. ,,Moc silnika'') wraz z zestawem
  etykiet (\code{LinguisticVariable}).
\item \textbf{$\Sigma$-count} --- suma stopni przynależności (uogólnienie ,,liczby
  elementów'' na zbiory rozmyte). To kluczowe narzędzie liczbowe całego projektu.
\end{itemize}

% #####################################################################
\part*{Część I --- MK (rozwinięta)}
\addcontentsline{toc}{section}{\textbf{CZĘŚĆ I --- MK}}

% =====================================================================
\section{Kardynalność (liczność) zbioru rozmytego}

\pyt{Kardynalność / liczba kardynalna zbioru rozmytego.}

W zbiorze \emph{klasycznym} kardynalność to po prostu liczba elementów. W zbiorze
\emph{rozmytym} element należy ,,po trochu'', więc zamiast liczyć sztuki, \textbf{sumujemy
stopnie przynależności}. To tzw.\ \textbf{$\Sigma$-count} (skalarna liczność):
\[
  \big|A\big| \;=\; \Sigma\text{count}(A) \;=\; \sum_{x\in X} \mu_A(x).
\]
Intuicja: jeśli trzy auta należą do zbioru ,,mocny'' w stopniach $1{,}0$, $0{,}6$ i $0{,}3$,
to ,,jest tam'' $1{,}9$ mocnego auta --- nie całe trzy, bo dwa pasują tylko częściowo.

\medskip
\textbf{Względna kardynalność} dzieli liczność przez rozmiar przestrzeni:
\(
  \|A\| = |A| / |X|.
\)
W projekcie liczona numerycznie jako \emph{średni stopień przynależności} po próbkach
przestrzeni --- \code{FuzzySet.relativeSigmaCardinality()}:
\[
  \|A\| \;=\; \frac{1}{|X|}\sum_{x\in X}\mu_A(x).
\]

\medskip
\textbf{Gdzie to działa w projekcie.} Cały aparat miar jakości stoi na $\Sigma$-countach:
stopień prawdziwości $T_1$ liczy $\frac{1}{m}\sum_i \mu_S(d_i)$ (to $\Sigma$-count
sumaryzatora po rekordach), a miary kardynalności $T_7,T_8,T_{10}$ używają względnej
kardynalności zbiorów ($\code{relativeSigmaCardinality()}$).

\medskip
\textbf{Uwaga --- dwie różne ,,licznoci''.} Projekt rozróżnia:
\begin{itemize}[leftmargin=*]
\item \emph{liczność $\Sigma$-count} = $\sum \mu(x)$ --- bierze pod uwagę kszta\l t funkcji
  (pole pod nią);
\item \emph{liczność nośnika} = ile punktów ma $\mu>0$ (\code{relativeSupportCardinality()})
  --- bierze tylko szerokość, ignoruje wysokość. To jest ,,poziom rozmytości'' $\inn(A)$
  (patrz §\ref{sec:fuzziness}).
\end{itemize}

% =====================================================================
\section{Suma i iloczyn zbiorów rozmytych (i alternatywne normy)}
\label{sec:union-inter}

\pyt{Suma, iloczyn zbioru --- co to, do czego, oraz ,,inaczej suma'' (s-norma).}

To rozmyte odpowiedniki logicznych \textbf{OR} i \textbf{AND}. W projekcie zdefiniowane
w \code{Operations} (dekoratory funkcji przynależności) i wystawione przez
\code{FuzzySet.union/intersect}:
\[
  \text{Iloczyn (AND):}\quad \mu_{A\cap B}(x) = \min\big(\mu_A(x),\,\mu_B(x)\big),
\qquad
  \text{Suma (OR):}\quad \mu_{A\cup B}(x) = \max\big(\mu_A(x),\,\mu_B(x)\big).
\]

\textbf{Do czego służą.} Budują \emph{złożone} sumaryzatory i kwalifikatory:
\begin{itemize}[leftmargin=*]
\item \emph{Iloczyn} --- ,,auto mocne \textbf{i} ciężkie''. Liczymy $\min$: o ocenie
  decyduje \textbf{najsłabsza} cecha (rygorystyczne ,,oba warunki naraz''). Auto mocne na
  $0{,}9$, ale lekkie (ciężkie na $0{,}2$) $\to$ iloczyn $=0{,}2$.
\item \emph{Suma} --- ,,auto szybkie \textbf{lub} tanie''. Liczymy $\max$: wystarczy jedna
  mocna cecha. Auto powolne ($0{,}1$), ale tanie ($0{,}8$) $\to$ suma $=0{,}8$.
\end{itemize}

\paragraph{,,Inaczej suma'' --- s-norma (i t-norma).} $\min$ i $\max$ to tylko
\emph{jeden, najpopularniejszy} wybór. Ogólnie:
\begin{itemize}[leftmargin=*]
\item iloczyn realizuje dowolna \textbf{t-norma} $T(a,b)$ (przemienna, łączna, monotoniczna,
  $T(a,1)=a$),
\item suma --- dowolna \textbf{s-norma} (t-konorma) $S(a,b)$ (jak wyżej, ale $S(a,0)=a$).
\end{itemize}
Najczęstsze pary (dualne przez dopełnienie $1-\cdot$):
\[
\begin{array}{lll}
\text{minimum / maksimum:} & T=\min(a,b) & S=\max(a,b)\\[2pt]
\text{algebraiczna:} & T=a\cdot b & S=a+b-a\cdot b\\[2pt]
\text{Łukasiewicza:} & T=\max(0,a+b-1) & S=\min(1,a+b)
\end{array}
\]
Czyli ,,suma algebraiczna'' $a+b-ab$ to alternatywna s-norma; nasz projekt świadomie używa
pary $\min/\max$ (komentarz w \code{Operations}: ,,t-norma minimum'', ,,t-konorma
maksimum''). $\max$ jest \emph{najmniejszą} możliwą s-normą, $\min$ --- \emph{największą}
t-normą.

% =====================================================================
\section{Dopełnienie zbioru}

\pyt{Dopełnienie zbioru rozmytego.}

Rozmyte \textbf{NOT}: ,,w jakim stopniu $x$ \emph{nie} pasuje do $A$''. Standardowo
(\code{Operations.complement}):
\[
  \mu_{\neg A}(x) \;=\; 1 - \mu_A(x).
\]
Przykład: $210$~KM należy do ,,mocny'' w stopniu $0{,}7$, więc do ,,nie mocny'' w stopniu
$0{,}3$. W kodzie \code{FuzzySet.complement()} tworzy nowy zbiór o nazwie ,,nie \dots''.

\medskip
\textbf{Ważna różnica względem zbiorów klasycznych.} Dla zbiorów rozmytych
\emph{nie zachodzą} prawa wyłączonego środka i sprzeczności:
\[
  A \cup \neg A \neq X, \qquad A \cap \neg A \neq \varnothing.
\]
Np.\ przy $\mu_A(x)=0{,}5$ mamy $\mu_{A\cap\neg A}(x)=\min(0{,}5;0{,}5)=0{,}5\neq 0$. To
nie błąd --- to istota rozmytości (element może być ,,trochę w'' i ,,trochę poza'').

% =====================================================================
\section{Operacje na zbiorach --- zestawienie i własności}

\pyt{Operacje na zbiorach.}

Trzy podstawowe (wszystkie w \code{Operations} / \code{FuzzySet}), liczone
\emph{punktowo} po przestrzeni:
\begin{center}
\begin{tabular}{lll}
\toprule
Operacja & Wzór & Logika \\
\midrule
Dopełnienie & $1-\mu_A(x)$ & NOT \\
Iloczyn & $\min(\mu_A(x),\mu_B(x))$ & AND (t-norma) \\
Suma & $\max(\mu_A(x),\mu_B(x))$ & OR (s-norma) \\
\bottomrule
\end{tabular}
\end{center}

\textbf{Własności, które warto znać:} przemienność, łączność, rozdzielność,
\emph{prawa De Morgana} ($\neg(A\cup B)=\neg A\cap\neg B$ --- zachodzą dla pary
$\min/\max$), idempotentność ($A\cap A=A$). \emph{Nie} zachodzą prawa
sprzeczności/wyłączonego środka (patrz wyżej). Warunek poprawności: oba zbiory muszą być
nad \textbf{tą samą przestrzenią} --- w kodzie \code{union}/\code{intersect} przyjmują
zbiory nad tym samym \code{Universe}.

% =====================================================================
\section{Miara T8 --- kardynalność sumaryzatora}

\pyt{Cechy miary T8 (stopień kardynalności sumaryzatora).}

\textbf{Za co odpowiada.} Mierzy, jak ,,ostry'' (skupiony) jest sumaryzator $S$ ---
podobnie jak nieprecyzyjność $T_2$, ale licząc na \emph{kardynalnościach} (pole pod
funkcją), a nie tylko na szerokości nośnika. Dzięki temu odróżnia trójkąt od trapezu o tej
samej podstawie. To miara \textbf{strukturalna} --- zależy od projektu zbiorów rozmytych,
\emph{nie} od danych w bazie.
\[
  T_8 \;=\; 1 - \sqrt[n]{\;\prod_{j=1}^{n} \frac{|S_j|}{|X_j|}\;}
       \;=\; 1 - \sqrt[n]{\;\prod_{j=1}^{n}\|S_j\|\;},
\]
gdzie $\|S_j\|$ to względna kardynalność $j$-tego atomu (w kodzie
\code{relativeSigmaCardinality()}), a $n$ --- liczba etykiet w sumaryzatorze
(średnia \emph{geometryczna} po atomach).

\textbf{Zakres i interpretacja.} $[0,1]$; \textbf{wyżej $=$ ostrzejszy} (węższy względem
uniwersum) sumaryzator. W projekcie typowo $\approx 0{,}88$--$0{,}95$ (etykiety są wąskie
względem pełnego zakresu cechy).

\textbf{Implementacja:} \code{QualityMeasures}, \code{t8 = 1 - geometricMean(cardS)}.

% =====================================================================
\section{Miara T11 --- długość kwalifikatora}

\pyt{Cechy miary T11.}

\textbf{Za co odpowiada.} Symetryczny odpowiednik $T_5$ (długość \emph{sumaryzatora}), ale
dla \textbf{kwalifikatora} $W$: krótszy kwalifikator (mniej etykiet) $=$ czytelniejsze
zdanie. Czysto \textbf{formalna} kara za złożoność --- nie zależy ani od danych, ani od
kształtu zbiorów.
\[
  T_{11} \;=\; 2\cdot(0{,}5)^{|W|},
\]
gdzie $|W|$ to liczba etykiet w kwalifikatorze.

\textbf{Zakres:} $(0,1]$; $|W|=1\Rightarrow T_{11}=1$, i maleje \emph{wykładniczo o połowę}
z każdą dodaną etykietą ($2\to0{,}5$, $3\to0{,}25$).

\textbf{Kluczowa cecha --- istnieje tylko w formie 2.} W formie 1 (bez kwalifikatora)
$T_{11}=0$ z definicji. \emph{Uwaga:} $T_{11}$ \textbf{nie występuje w książce źródłowej}
(Niewiadomski definiuje $T_1$--$T_{10}$). W projekcie dodana przez analogię do $T_5$, żeby
kwalifikator też miał miarę długości. Implementacja: \code{t11 = 2.0 * Math.pow(0.5, mr)}.

% =====================================================================
\section{Dwie formy wyrażeń kwantyfikowanych lingwistycznie}
\label{sec:forms}

\pyt{Dwie formy wyrażeń kwantyfikowanych lingwistycznie (formy 1 i 2 podsumowania).}

To dwa szablony zdań Zadeha/Yagera. W projekcie buduje je \code{SingleSubjectSummary}.

\paragraph{Forma 1 (protozdanie, bez kwalifikatora): ,,$Q\ P$ są $S$''.}
\[
  T_1 = \mu_Q\!\Big(\tfrac{1}{m}\textstyle\sum_{i=1}^m \mu_S(d_i)\Big)
  \quad(\text{kwant. względny}),\qquad
  T_1 = \mu_Q\!\big(\textstyle\sum_i \mu_S(d_i)\big)\ (\text{bezwzględny}).
\]
Przykład: \emph{,,Większość aut ma standardową długość''}. Liczymy proporcję aut
spełniających $S$ i sprawdzamy, czy trafia w kwantyfikator $Q$.

\paragraph{Forma 2 (z kwalifikatorem $W$): ,,$Q\ P$ będących $W$ są $S$''.}
W kodzie: \emph{,,[Q] aut, które są [W], ma [S]''}. Kwalifikator $W$ zawęża populację
(rozmyty odpowiednik \code{WHERE}). Stopień prawdziwości liczony \textbf{warunkowo}:
\[
  T_1 = \mu_Q\!\left(\frac{\sum_{i=1}^m \min\big(\mu_W(d_i),\,\mu_S(d_i)\big)}
                          {\sum_{i=1}^m \mu_W(d_i)}\right).
\]
Przykład: \emph{,,Większość ciężkich aut ma dużą moc''} ($W=$ ,,ciężki'').

\textbf{Różnica praktyczna.} Forma 1 mówi o \emph{całej} bazie; forma 2 --- o
\emph{wycinku} (autach spełniających $W$). Forma 2 aktywuje dodatkowe miary $T_9,T_{10},
T_{11}$ (jakość kwalifikatora), których w formie 1 nie ma (są zerowe).

% =====================================================================
\section{Zadanie: policzenie kardynalności dla zbioru}

\pyt{Policzyć kardynalność dla (rozmytego) zbioru.}

Schemat liczenia (przykład reprezentatywny --- liczby zależą od konkretnego zadania).
Niech przestrzeń to pięć aut, a zbiór rozmyty $A=$,,mocny silnik'' ma stopnie:
\[
  A = \frac{0{,}2}{a_1} + \frac{0{,}8}{a_2} + \frac{1{,}0}{a_3} + \frac{0{,}5}{a_4}
      + \frac{0{,}0}{a_5}.
\]

\begin{enumerate}[leftmargin=*]
\item \textbf{Kardynalność ($\Sigma$-count)} --- suma stopni:
  \[
    |A| = 0{,}2+0{,}8+1{,}0+0{,}5+0{,}0 = \boxed{2{,}5}.
  \]
\item \textbf{Względna kardynalność} --- podziel przez $|X|=5$:
  \[
    \|A\| = 2{,}5/5 = \boxed{0{,}5}.
  \]
\item \textbf{Nośnik i jego liczność} --- elementy z $\mu>0$:
  $\text{supp}(A)=\{a_1,a_2,a_3,a_4\}$, więc $|\text{supp}(A)|=4$, a
  \emph{poziom rozmytości} $\inn(A)=4/5=0{,}8$.
\item \textbf{Wysokość} $=\max_x\mu_A(x)=1{,}0$ $\Rightarrow$ zbiór \textbf{normalny}.
\end{enumerate}

\noindent\emph{Kontrast klasyczny:} gdyby zbiór był ostry (tylko $0$/$1$), kardynalność
$=$ liczba elementów z przynależnością $1$. $\Sigma$-count to gładkie uogólnienie: zlicza
,,częściowe'' przynależności.

% #####################################################################
\part*{Część II --- AN (skrótowo)}
\addcontentsline{toc}{section}{\textbf{CZĘŚĆ II --- AN}}

% =====================================================================
\section{Zbiory, nośnik, $\alpha$-przekrój, poziom rozmytości}

\pytan{Nośnik zbioru i $\alpha$-przekrój.}
\textbf{Nośnik} $\text{supp}(A)=\{x:\mu_A(x)>0\}$ --- wszystko, co należy choć trochę.
\textbf{$\alpha$-przekrój} $A_\alpha=\{x:\mu_A(x)\ge\alpha\}$ --- wszystko, co należy
\emph{przynajmniej w stopniu $\alpha$}; ostry (klasyczny) zbiór ,,wycięty'' na wysokości
$\alpha$. Im wyższe $\alpha$, tym węższy przekrój. Nośnik $\approx$ przekrój dla $\alpha\to
0^+$, a rdzeń (core) $=A_1$. W kodzie: \code{FuzzySet.support()} i \code{alphaCut(alpha)}
(zwracają przedział obejmujący).

\pytan{Co to i do czego ,,poziom (stopień) rozmytości''.}
\label{sec:fuzziness}
To miara, jak ,,rozlany'' jest zbiór: $\inn(A)=|\text{supp}(A)|/|X|$ (w kodzie
\code{relativeSupportCardinality()}). $\inn(A)$ blisko $0$ $=$ wąska, precyzyjna etykieta;
blisko $1$ $=$ rozmazana na całe uniwersum. \textbf{Do czego służy:} jest cegiełką miar
nieprecyzyjności --- $T_2=1-\sqrt[n]{\prod \inn(S_j)}$, $T_6=1-\inn(Q)$, $T_9=1-\inn(W)$.

\pytan{Zbiór rozmyty i jego zastosowanie w reprezentacji informacji lingwistycznych.}
Zbiór rozmyty modeluje \emph{nieostre pojęcie języka naturalnego} (,,mocny'', ,,duży'',
,,szybki''), które nie ma twardej granicy. Funkcja $\mu$ tłumaczy słowo na liczbę: pozwala
maszynie ocenić ,,$190$~KM jest mocne w stopniu $0{,}7$''. Dzięki temu możemy budować i
\emph{oceniać} zdania w quasi-naturalnym języku --- to fundament całego projektu
(\code{FuzzySet} $\to$ \code{Label} $\to$ \code{LinguisticVariable} $\to$ podsumowania).

% =====================================================================
\section{Zmienna lingwistyczna}

\pytan{Zmienna lingwistyczna i jej zastosowanie.}
To \textbf{cecha} (np.\ ,,Moc silnika'') wraz z: przestrzenią rozważań, kolumną bazy i
\textbf{zestawem etykiet} (słów: ,,słaby'', ,,średni'', ,,mocny''), z których każda jest
osobnym zbiorem rozmytym (\code{LinguisticVariable} + lista \code{Label}). \textbf{Zastosowanie:}
zamienia surową liczbę z bazy (np.\ $130$~KM) na słowa ze stopniami przynależności --- to
warstwa ,,liczby $\to$ słowa'', bez której nie da się ułożyć zdania podsumowania.

% =====================================================================
\section{Kwantyfikatory}

\pytan{Kwantyfikatory względne i bezwzględne --- zastosowanie.}
\textbf{Względny} działa na \emph{proporcji} $[0,1]$ (,,większość'', ,,około połowy'') ---
mówi, jaka \emph{część} obiektów ma cechę. \textbf{Bezwzględny} działa na \emph{liczbie}
sztuk (,,około 100 aut'') --- mówi o liczności. W projekcie (\code{Quantifier.Type})
względne mają trapezy na $[0,1]$, np.\ ,,większość'' $=(0{,}60;0{,}70;0{,}85;0{,}95)$,
,,około połowy'' $=(0{,}30;0{,}40;0{,}60;0{,}70)$. Zastosowanie: argumentem $\mu_Q$ jest
proporcja (względny) lub sama suma $\Sigma$-count (bezwzględny) --- patrz $T_1$ w §\ref{sec:forms}.

\pytan{Kwantyfikator rozmyty a kwantyfikator lingwistyczny --- różnica.}
\textbf{Lingwistyczny} to samo \emph{słowo} języka naturalnego (,,większość'') ---
pojęcie. \textbf{Rozmyty} to jego \emph{matematyczny model}: zbiór rozmyty z funkcją
przynależności, którą da się policzyć. Innymi słowy: kwantyfikator lingwistyczny
\emph{reprezentujemy} kwantyfikatorem rozmytym. W kodzie świetnie to widać: \code{Quantifier}
trzyma \code{name} (warstwa lingwistyczna) \emph{oraz} \code{FuzzySet set} (warstwa rozmyta,
czyli model).

% =====================================================================
\section{Paradoks Poincaré'a}

\pytan{Paradoks Poincaré'a.}
Dotyczy \textbf{nierozróżnialności}. Masy $a$ i $a{+}1$~grama są dla zmysłów
nieodróżnialne ($a\approx a{+}1$), tak samo $a{+}1\approx a{+}2$ itd., ale $a$ i
$a{+}1000$ \emph{są} rozróżnialne. Zatem relacja ,,nierozróżnialne'' jest \textbf{zwrotna}
i \textbf{symetryczna}, lecz \textbf{nie przechodnia} --- łańcuch ,,równych'' kroków daje
różnicę widoczną. To pokazuje, że klasyczna relacja równoważności (wymagająca
przechodniości) \emph{nie nadaje się} do modelowania percepcji/podobieństwa --- stąd
potrzeba \textbf{relacji sąsiedztwa} (tolerancji) i logiki rozmytej. Łączy się to wprost z
pytaniami o relacje niżej.

% =====================================================================
\section{Relacje rozmyte: zwrotność, symetria, przechodniość}

\pytan{Zwrotność i symetria relacji rozmytej --- przykłady z języka naturalnego.}
Relacja rozmyta to $R:X\times X\to[0,1]$ ($R(x,y)=$ stopień, w jakim $x$ jest w relacji z $y$).
\begin{itemize}[leftmargin=*]
\item \textbf{Zwrotność:} $R(x,x)=1$. NL: ,,każde słowo jest w 100\% podobne do samego
  siebie''; ,,auto jest tak samo szybkie jak ono samo''.
\item \textbf{Symetria:} $R(x,y)=R(y,x)$. NL: ,,jeśli Kowalski jest podobny do Nowaka, to
  Nowak jest tak samo podobny do Kowalskiego''; ,,bliski'', ,,sąsiad''.
\end{itemize}

\pytan{Relacja rozmyta przechodnia (tranzytywna) --- definicja.}
\textbf{Przechodniość typu max-min:}
\[
  R(x,z)\;\ge\;\max_{y}\ \min\big(R(x,y),\,R(y,z)\big)\qquad\forall x,z.
\]
NL: jeśli $x$ jest podobny do $y$, a $y$ do $z$, to $x$ musi być podobny do $z$
\emph{przynajmniej} w stopniu najsłabszego ogniwa po najlepszej drodze. (Ogólnie zamiast
$\min$ dowolna t-norma --- $T$-przechodniość.)

% =====================================================================
\section{Podobieństwo: sąsiedztwo vs rozmyta równoważność}

\pytan{Relacja sąsiedztwa a relacja podobieństwa --- cechy i definicja.}
\pytan{Podobieństwo jako relacja sąsiedztwa czy jako rozmyta równoważność?}
\begin{itemize}[leftmargin=*]
\item \textbf{Relacja sąsiedztwa} (tolerancji, proximity): \textbf{zwrotna + symetryczna},
  \emph{bez} wymogu przechodniości. Rozmyty odpowiednik relacji tolerancji.
\item \textbf{Relacja podobieństwa} w sensie Zadeha $=$ \textbf{rozmyta równoważność}:
  zwrotna + symetryczna \textbf{+ przechodnia} (max-min). Rozmyty odpowiednik relacji
  równoważności.
\end{itemize}
\textbf{Cała różnica to przechodniość.} I tu wraca Poincaré: realne ,,podobieństwo''
percepcyjne \emph{łamie} przechodniość (łańcuch podobnych obiektów kończy się parą
niepodobnych), więc naturalniej modelować je jako \textbf{relację sąsiedztwa}, a nie jako
ścisłą rozmytą równoważność. Równoważność rozmyta to silniejsze, idealizowane założenie
(pozwala na rozmyte klasy abstrakcji).

\pytan{nGram --- czy jest symetryczna?}
Miara podobieństwa oparta o $n$-gramy (liczba/udział wspólnych $n$-gramów dwóch napisów,
np.\ współczynnik Dice'a) jest \textbf{symetryczna} --- ,,wspólne $n$-gramy $a$ i $b$'' to
to samo co ,,wspólne $n$-gramy $b$ i $a$'', więc $\text{sim}(a,b)=\text{sim}(b,a)$. Jest
też zwrotna ($\text{sim}(a,a)=1$), ale \textbf{nie} przechodnia --- więc jest
\emph{relacją sąsiedztwa}, a nie rozmytą równoważnością.

\pytan{Miara podobieństwa a metryka --- różnice.}
\begin{center}
\begin{tabular}{lll}
\toprule
 & Metryka $d(x,y)$ & Miara podobieństwa $s(x,y)$ \\
\midrule
Kierunek & $0=$ identyczne, rośnie z różnicą & $1=$ identyczne, maleje z różnicą \\
Tożsamość & $d(x,x)=0$ & $s(x,x)=1$ (zwrotność) \\
Symetria & tak & zwykle tak \\
Nierówność trójkąta & \textbf{wymagana} & \textbf{nie} wymagana \\
Zakres & $[0,\infty)$ & zwykle $[0,1]$ \\
\bottomrule
\end{tabular}
\end{center}
Krótko: metryka mierzy \emph{odległość} i musi spełniać nierówność trójkąta; miara
podobieństwa mierzy \emph{bliskość} (odwrotnie) i nie musi. Często $s=1-d_{\text{znorm.}}$,
ale nie każde podobieństwo pochodzi od metryki (brak nierówności trójkąta to dokładnie
sytuacja sąsiedztwa/Poincaré).

% =====================================================================
\section{Miary jakości --- pojęcia AN}

\pytan{Znaczenie miary ,,stopień nieprecyzyjności''.}
Mówi, jak \textbf{ostre} są użyte etykiety (a nie co jest w danych). $T_2=1-\sqrt[n]{\prod
\inn(S_j)}$ dla sumaryzatora (analogicznie $T_6$ dla kwantyfikatora, $T_9$ dla
kwalifikatora). \textbf{Wyżej $=$ precyzyjniej} (wąskie zbiory względem uniwersum). To
miara \emph{strukturalna} --- premiuje zdania zbudowane z ostrych pojęć, niezależnie od
zawartości bazy.

\pytan{Miara T3 (stopień pokrycia) --- co oznacza dla podsumowania.}
Mówi, \textbf{jak dużą część rekordów} podsumowanie faktycznie obejmuje (udział aut z
$\mu_S>0$; w formie 2 --- udział wśród spełniających kwalifikator). $T_3=0{,}694$ znaczy
,,dotyczy 69{,}4\% aut''. Razem z $T_1$ to jedna z dwóch miar realnie zależnych od danych:
$T_1$ mówi \emph{czy prawda}, $T_3$ --- \emph{o ilu obiektach}. Niskie $T_3$ nie zawsze
jest złe (rzadkie, ale prawdziwe i nieoczywiste zdanie).

\pytan{Cechy miary T5 (dwie).}
$T_5=2\cdot(0{,}5)^{|S|}$ (długość sumaryzatora). Dwie cechy:
\begin{enumerate}[leftmargin=*]
\item \textbf{Zależy wyłącznie od liczby etykiet} $|S|$ --- nie od danych, nie od kształtu
  ani szerokości zbiorów, nie od kwantyfikatora. Czysto formalna kara za złożoność
  (czytelność).
\item \textbf{Maleje wykładniczo o połowę}: $|S|{=}1\to1$ (najkrótsze, max), $2\to0{,}5$,
  $3\to0{,}25$; wartości w $(0,1]$.
\end{enumerate}

% =====================================================================
\section{Forma 2 podsumowania --- i co liczy się szybciej}

\pytan{Druga forma jednopodmiotowego podsumowania lingwistycznego.}
,,$Q$ aut, które są $W$ (kwalifikator), ma $S$ (sumaryzator)'' --- z kwalifikatorem
zawężającym populację. $T_1$ liczony warunkowo (wzór w §\ref{sec:forms}).

\pytan{Co liczy się szybciej: $S_1$ i $S_2$ bez kwalifikatora (forma 1, sumaryzator
złożony) czy $S_1$ z $Q$ i $S_2$ jako kwalifikator (forma 2)?}
Obie wersje robią \textbf{jeden przebieg} po rekordach, więc są tej samej klasy złożoności
$O(m\cdot\#\text{atomów})$. Ale \textbf{forma 1 (z $S_2$ wbudowanym w sumaryzator} jako
$\min(\mu_{S_1},\mu_{S_2})$) liczy się \textbf{taniej}, bo:
\begin{itemize}[leftmargin=*]
\item potrzebuje tylko jednego akumulatora ($\sum\mu_S$) i prostej średniej, bez
  warunkowego dzielenia przez $\sum\mu_W$;
\item \emph{nie} uruchamia miar kwalifikatora $T_9,T_{10},T_{11}$ (w formie 1 są zerowe).
\end{itemize}
W kodzie widać to wprost: gałąź \code{form2} w \code{QualityMeasures} dokłada akumulatory
\code{sumR}, \code{sumRandS}, \code{countRpos}, \code{countRSpos} oraz cały blok
$T_9$--$T_{11}$. Wniosek: \textbf{przeniesienie $S_2$ do sumaryzatora (forma 1) jest
szybsze} niż użycie go jako osobnego kwalifikatora (forma 2), choć rząd złożoności ten sam.
Semantycznie jednak liczą \emph{co innego} (forma 1 --- ,,$Q$ aut jest $S_1$ \emph{i}
$S_2$''; forma 2 --- ,,$Q$ \emph{spośród} aut $S_2$ jest $S_1$'').

% =====================================================================
\section*{Szybka ściąga --- jednym zdaniem}
\addcontentsline{toc}{section}{Szybka ściąga}
\begin{itemize}[leftmargin=*,nosep]
\item Iloczyn $=\min$ (t-norma, AND), suma $=\max$ (s-norma, OR), dopełnienie $=1-\mu$.
\item Kardynalność $=\Sigma$-count $=\sum\mu$; względna $=$ to / $|X|$.
\item Poziom rozmytości $\inn(A)=|\text{supp}|/|X|$; nieprecyzyjność $=1-\inn$.
\item Nośnik: $\mu>0$; $\alpha$-przekrój: $\mu\ge\alpha$.
\item Sąsiedztwo $=$ zwrotna+symetryczna; podobieństwo (rozmyta równoważność) $=$ jeszcze
  przechodnia. Poincaré: percepcja łamie przechodniość $\Rightarrow$ sąsiedztwo.
\item nGram: symetryczna i zwrotna, nieprzechodnia.
\item Metryka: $0=$ identyczność, nierówność trójkąta wymagana; podobieństwo: $1=$
  identyczność, bez trójkąta.
\item $T_3$ pokrycie + $T_1$ prawda $=$ jedyne miary ,,o danych''; reszta strukturalna/formalna.
\end{itemize}

\end{document}
